\begin{lemma}\label{lem:avoid}
The map \(R\) of Definition~\ref{def:blaschke-quotient} has no zero and no pole on the closed annulus \(1\le|z|\le r'\). In particular \(u(z)=\log|R(z)|\) is harmonic on an open neighborhood of \(\overline{\mathbb A(1,r')}\).
\end{lemma}

\begin{proof}
The zeros and poles of a Blaschke factor and its reciprocal occur among \(a,\ b,\ 1/\overline a,\ 1/\overline b\). The points \(a\in\mathcal A\), \(b\in\mathcal B\) lie in the unit disk, with
\[ \max_{a\in\mathcal A}|a|=\frac{\sqrt{53}}{10},\qquad \max_{b\in\mathcal B}|b|=\frac{\sqrt{40}}{10}. \]
Direct squaring gives
\[ \Bigl(\tfrac{1287}{1000}\Bigr)^2\tfrac{53}{100}<1,\qquad \Bigl(\tfrac{1287}{1000}\Bigr)^2\tfrac{40}{100}<1, \]
so every reciprocal point \(1/\overline a\), \(1/\overline b\) lies outside \(|z|=r'\). Hence \(R\) has neither zero nor pole on \(1\le|z|\le r'\), and \(u=\log|R|\) is harmonic on a neighborhood of the closed annulus.
\end{proof}